%% ============================================================
%% deliverables.tex — 산출물 (본문, 참고문헌 앞)
%% 보고 기간에 작성한 주요 산출물의 요약. 전문은 참고문헌 뒤 부록에 수록.
%% ============================================================

\chaptersummary{보고 기간에 산출한 8종의 결과물을 요약하고 그 성격을 정리한다. 소프트웨어 문서화·추적성·형상관리 관점에서 각 산출물의 위상을 밝힌다.}
\chapter{산출물 개요}
보고 기간 동안 ARES 프로젝트에서 작성한 주요 문서$\cdot$코드 산출물의 요약이다.
각 산출물의 \textbf{전문은 참고문헌 뒤 부록}(부록 A$\sim$I)에 수록한다. 앞의 세 산출물
(부록 A$\sim$C)이 \emph{기준선 분석과 개선 설계}라면, 부록 D$\cdot$E는 그 설계가 통합
머지로 실행된 뒤의 \emph{중간 점검$\cdot$재분석}이고, 부록 F$\cdot$G는 12일차의
\emph{품질 실측(테스트 결과)과 도구 사용법}이며, 부록 H는 14일차 \emph{최종 산출물
코드}를 부록 A 구조로 정리하고 \emph{변화와 개선점}을 종합한 것이고, 부록 I는 그
\emph{최종 산출물 소스 코드(펌웨어$\cdot$웹 코어)를 원문 그대로} 수록한 것이다.

\begin{center}
\renewcommand{\arraystretch}{1.3}
\begin{tabularx}{\linewidth}{@{}>{\sffamily\bfseries}l L{0.46\linewidth} >{\sffamily}c@{}}
\toprule
산출물 & 내용 & 전문 \\
\midrule
코드 분석 종합 보고서 & 저장소 전체의 코드 수준 기술 분석(전 14장) & 부록 A \\
1차 Refactoring 계획 & 비대 파일$\cdot$함수의 동작 보존 구조 분할 설계 & 부록 B \\
2차 Refactoring 계획 & 시뮬$\cdot$피코$\cdot$웹 3트랙 고도화 계획 & 부록 C \\
중간 점검 머지 --- 주요 코드 변경 & 통합 머지로 반영된 부록 A 대비 변경의 트랙별 요약 & 부록 D \\
중간 보고 시점 코드 분석 & 부록 A 장 구조를 현재 코드에 재적용한 변화 중심 재분석 & 부록 E \\
소프트웨어 테스트 결과 보고서(STR v1.1) & TCL v1.0(98케이스) 1차 실측 --- 85 Pass$\cdot$0 Fail, 조건부 Go & 부록 F \\
시뮬레이터 사용 매뉴얼 & 씬 생성 도구(컴포넌트 9종$\cdot$개발자/사용자 모드) 사용 안내 --- 실구동 캡처 수록 & 부록 G \\
최종 산출물 코드 정리 & 14일차 최종 코드를 부록 A 구조로 재정리 --- 트랙별 변화와 개선점 종합 & 부록 H \\
최종 산출물 소스 코드 리스팅 & 펌웨어(\code{Pico/})$\cdot$웹 코어(\code{Web/}) 소스 원문 --- 저장소와 바이트 일치 & 부록 I \\
\bottomrule
\end{tabularx}
\end{center}

\section{코드 분석 종합 보고서 (요약)}
ARES 저장소~\cite{ares-repo} 전체를 코드 수준에서 분석한 \textbf{「ARES 프로젝트
코드 분석 종합 보고서」(전 14장)}이다. ARES는 세 런타임 도메인 — 웹 UI(Blockly +
Web Bluetooth), 피코 펌웨어(MicroPython), AI 보조 모듈 — 으로 구성되며, 이들은
\emph{개행 구분 텍스트 명령 프로토콜}이라는 단일 계약으로 느슨하게 결합된다.

\begin{itemize}
  \item \textbf{느슨한 결합} — 웹과 펌웨어는 코드를 공유하지 않고 텍스트 명령만으로
  통신해 독립 개발$\cdot$교체가 쉽다. 반면 프로토콜 변경 시 양측 동시 반영 의존이 생기며,
  컴파일 검증이 없어 \code{API\_DOCUMENTATION.md}가 사실상의 계약 문서다.
  \item \textbf{모듈 패턴} — 모든 하드웨어 모듈은 활성화 확인 → try/except 임포트 →
  실패 시 \code{None}의 동일 패턴을 따른다.
  \item \textbf{주요 제약} — MicroPython RAM 제약, BLE 20바이트 청크~\cite{bluetooth54}, 기기 이름 10자,
  Web Bluetooth의 HTTPS/localhost 요구.
\end{itemize}
\noindent 전문: \ref{app:code}.

\section{1차 Refactoring 계획 (요약)}
지나치게 큰 함수$\cdot$파일을 기능 단위로 분할하기 위한 \emph{동작 보존} 설계로,
비대 지점을 네 유형으로 분류했다~\cite{fowler2018refactoring, martin2008cleancode, feathers2004legacy, gof1994, mcconnell2004codecomplete, hunt1999pragmatic}.

\smallskip
\noindent\chip{A · 거대 클로저} \chip{B · 거대 디스패처} \chip{C · 데이터 덩어리} \chip{D · 인라인 자산}

\begin{itemize}
  \item \textbf{유형 A} — \code{simulation.js}(2,333줄)를 공유 컨텍스트(ctx) +
  서브시스템 14개로.
  \item \textbf{유형 B} — \code{process\_data.py}를 도메인 핸들러 + dict 디스패치로.
  \item \textbf{유형 C} — \code{blocklyconfig.js}를 카테고리별 파일로.
  \item \textbf{유형 D} — \code{dashboard.html}의 인라인 CSS$\cdot$JS를 외부 파일로.
\end{itemize}
\noindent 전문: \ref{app:refactor1}.

\section{2차 Refactoring 계획 (요약)}
1차 구조 분할의 연장선에서 \textbf{시뮬레이터 $\cdot$ 피코+블록코딩 $\cdot$ 웹 인터페이스}
세 트랙을 함께 끌어올린다. 시뮬레이터 요소별 분할과 Transform 도구화, 미션별 블록
노출, 반응형$\cdot$모바일 최적화$\cdot$디자인 현대화를 다룬다.
\noindent 전문: \ref{app:refactor2}.

\section{중간 점검 머지 --- 주요 코드 변경 (요약)}
1$\cdot$2차 Refactoring 계획이 \textbf{작업 브랜치 통합 머지}로 실제 코드에 반영된 결과를,
부록 A(기준선)와 대조해 트랙별로 정리한 \emph{중간 점검}이다. 코드 도메인에서
\textbf{35개 파일 +5{,}152 / $-$2{,}563 라인}이 바뀌었고, 핵심은 유형 A(거대 클로저)
분할로 \code{simulation.js}($\sim$2{,}300$\to$383줄)를 \emph{공유 컨텍스트 + 14 서브시스템}
으로 나눈 것, 피코 설정을 \code{system.json}$\to$\code{ModelFactorySetting.txt}로 텍스트화한
것, 명령 프로토콜에 속도 인자$\cdot$\code{GET\_NAMES}를 확장한 것이다.
\noindent 전문: \ref{app:merge}.

\section{중간 보고 시점 코드 분석 (요약)}
부록 A의 장 구성(하드웨어$\cdot$핀 $\to$ 펌웨어 $\to$ 웹 계층 $\to$ 블루투스 $\to$ 블록코딩
$\to$ WebGL $\to$ 연계 구조 $\to$ API)을 \textbf{중간 보고 시점의 현재 코드에 다시 적용}해,
각 장에서 현재 구현과 \emph{[부록 A 대비] 변화}를 실제 코드 발췌로 드러낸 재분석이다.
부록 D가 변경의 \emph{요약}이라면 본 산출물은 부록 A 틀의 \emph{재분석}으로 상보적이며,
하드웨어 계층 무변동, 펌웨어의 중단 가능 sleep$\cdot$설정 텍스트화, 웹의 \code{ui.js}
분리$\cdot$BLE 응답 동기화, 시뮬레이터의 dt 기반 애니메이션$\cdot$\code{Web/Simulation/}
죽은 코드 확인 등을 담는다.
\noindent 전문: \ref{app:midcode}.

\section{소프트웨어 테스트 결과 보고서 --- STR v1.1 (요약)}
11일차 QA 체계$\cdot$TCL v1.0(14 카테고리$\cdot$98 케이스)에 따라 12일차에 본격 수행한
\textbf{1차 소프트웨어 테스트의 결과 보고서}다~\cite{iso29119, ieee829, iso25010, beck2003tdd}. 실기(BLE)와 시뮬레이터를 병행해
\textbf{85건 Pass $\cdot$ 0건 Fail(86.7\%)}을 기록했고, 실행 완료 케이스에서 기능 결함은
발견되지 않았다. 서보$\cdot$DC$\cdot$OLED$\cdot$센서$\cdot$시뮬레이터 3D(24건)$\cdot$프리셋
E2E는 전건 Pass. 미결 13건은 전건 \emph{비결함}(테스트 인프라 5$\cdot$장비 4$\cdot$명세 2$\cdot$%
레거시 2)으로 분석되어 개선 항목 IMP-001$\sim$004(최우선: 실기 BLE 수신 로그 표시)로
관리하며, 판정은 \textbf{조건부 Go}(재시험 R1$\cdot$R2 완료 시 확정)다.
\noindent 전문: \ref{app:str} (원문 \code{Test/3\_STR.html}).

\section{시뮬레이터 사용 매뉴얼 (요약)}
12$\sim$13일차에 구축$\cdot$고도화한 \textbf{블록코딩 연계 씬 생성 도구}의 사용
매뉴얼로, 실제 서비스를 구동$\cdot$조작하며 캡처한 \textbf{실구동 화면 9장}을 수록했다.
사용자는 씬 드롭다운(기본 주제 + 서비스 등록 씬)에서 씬을 골라
\textbf{모의실행/실험중단}으로 블록 코드를 장치 없이 실행하고, 개발자는 \textbf{Ctrl+E}로
진입해 우클릭 스폰(기본 도형$\cdot$OLED$\cdot$GLB 자산 목록)$\cdot$복제(Ctrl+V)$\cdot$기본색/발광색
편집, 기즈모$\cdot$Hierarchy(길게 클릭 이름 변경)$\cdot$인스펙터(트랜스폼$\cdot$구조화 컴포넌트
필드$\cdot$사용/미사용 체크박스$\cdot$회전축 핸들 편집)로 씬을 구성한 뒤 JSON 저장 또는
\textbf{서비스 등록}(manifest 자동 갱신)으로 배포한다. \textbf{컴포넌트 9종 레퍼런스}
(LED$\cdot$Buzzer$\cdot$Oled$\cdot$DC$\cdot$Servo$\cdot$UltraSonic$\cdot$Magnet$\cdot$Metal$\cdot$Gun --- Gun
자기 비행$\cdot$복귀, LED 포인트 라이트, Servo neutral 포함)와 좌표계 규약(\textbf{축
방향=부모 좌표계$\cdot$기준점/감지점=객체 로컬}$\cdot$오프셋 m$\cdot$스케일 비전파),
작업 씬 임시 보존을 담았다.
\noindent 전문: \ref{app:simguide}.

\section{최종 산출물 코드 정리 (요약)}
14일차 \textbf{최종 산출물 코드}를 부록 A의 장 구조(하드웨어$\cdot$핀 $\to$ 펌웨어 $\to$ 웹
서비스 $\to$ 블루투스 $\to$ 블록코딩 $\to$ WebGL 시뮬레이션 $\to$ 연계 구조 $\to$ API)로 다시
정리한 분석이다. 하드웨어$\cdot$핀 계층은 기준선(부록 A) 그대로 안정적이고, 상위 계층은
\textbf{펌웨어}(중단 가능 sleep$\cdot$속도 클램프$\cdot$텍스트 공장 설정$\cdot$\code{GUN\_FIRE}
\code{fire\_once}), \textbf{웹}(\code{ui.js} 분리$\cdot$맥락 인지형 중앙 버튼$\cdot$동적 타임아웃),
\textbf{시뮬레이터}(14$\to$17 모듈, 씬 생성 도구$\cdot$색상/발광 시스템$\cdot$회전축 로컬 규약$\cdot$%
씬 서비스 등록$\cdot$그림자$\cdot$LED 포인트 라이트$\cdot$Gun 자기비행/복귀)로 크게 진전했다.
말미에 기준선(부록 A)$\cdot$중간(부록 E) 대비 \textbf{변화와 개선점}을 트랙별 표와 함께
종합한다.
\noindent 전문: \ref{app:finalcode}.


%% ===== 장 심화(관련 연구·의사코드·의의) — 자동 보강 =====
\section{관련 연구 및 배경}

본 장은 보고 기간에 산출된 문서와 코드 산출물(work product)을 정리하는 메타 장으로, 조사·의사코드·의의 모두 \emph{소프트웨어 산출물 관리·문서화·요구사항 추적성·품질}이라는 주제로 다룬다. 보고 기간의 산출물은 코드분석 종합보고서, 1차·2차 리팩토링 계획, 통합 머지 점검 리뷰, 중간 재분석 보고서, 소프트웨어 시험 보고서(STR), 시뮬레이터 사용 매뉴얼, 최종 코드정리 문서, 그리고 전체 소스 리스팅으로 구성된다. 이들은 각각 개발 생애주기의 서로 다른 단계를 문서화하며, 하나의 정합된 산출물 체계로 묶일 때 프로젝트의 재현성과 감사(audit) 가능성이 확보된다~\cite{mcconnell2004codecomplete,iso25010}.

소프트웨어 문서화의 학술적 근간은 산출물을 개발 활동의 부산물이 아니라 명시적 산출물로 관리하는 데 있다. 코드 자체가 진실의 근원이지만, 코드분석 보고서와 리팩토링 계획은 왜 그런 구조를 택했는지에 대한 설계 의도를 보존한다~\cite{martin2008cleancode,hunt1999pragmatic}. 특히 레거시 코드를 다루는 이 프로젝트에서, 기존 모놀리식 시뮬레이터를 모듈 라이브러리로 분해한 이력을 문서로 남기는 것은 이후 유지보수자가 변경의 맥락을 이해하도록 돕는다~\cite{feathers2004legacy,fowler2018refactoring}. 리팩토링 계획을 1차·2차로 나눈 것 자체가 대규모 구조 변경을 점진적·검증 가능한 단계로 쪼개는 실무 원칙을 반영한다.

시험 문서화 측면에서 STR(Software Test Report)은 국제 표준이 규정하는 시험 산출물 체계와 정렬된다. ISO/IEC/IEEE 29119는 시험 과정과 문서를 정의하고, IEEE 829는 시험 계획·설계·사례·절차·보고의 문서군을 규정하며, ISO/IEC 25010은 기능 적합성·신뢰성·사용성 등 품질 특성의 평가 틀을 제공한다~\cite{iso29119,ieee829,iso25010}. ARES의 STR은 BLE 텍스트 프로토콜 명령 처리, 시뮬레이터 명령 디스패치, 하드웨어 모듈 활성화 로직 등을 대상으로 기대 결과와 실제 결과를 대조하여 품질을 근거 있게 주장한다. 이러한 표준 정렬은 교육용 프로젝트라 하더라도 산출물이 전문적 품질 관행을 따르도록 만든다.

요구사항 추적성(traceability)은 산출물들을 가로지르는 연결 고리다. 미션 요구사항(예: 로버 이동, LED 제어, 거리·자기 센서 읽기)이 블록 정의, 명령 실행기, 펌웨어 명령 처리기, 시뮬레이터 효과, 그리고 STR의 시험 사례로 이어지는 사슬을 추적할 수 있어야 한다~\cite{iso29119,mcconnell2004codecomplete}. 추적성은 변경 영향 분석을 가능하게 하여, 예컨대 명령 프로토콜의 한 항목이 바뀔 때 영향받는 웹 UI·펌웨어·시뮬레이터·시험을 식별하도록 돕는다. 중간 재분석 보고서는 바로 이 추적 관계를 리팩토링 이후 시점에 다시 검증하는 산출물이다.

형상관리(configuration management)와 산출물 관리는 이 모든 문서의 버전과 상태를 통제한다. 프로젝트는 Git을 정본 저장소로 삼고, 대형 빌드 산출물(진행보고 PDF)은 릴리스 자산으로 분리하며, 웹 배포는 자동화 파이프라인으로 처리한다~\cite{chacon2014progit,ares-repo}. 이는 산출물마다 저장 위치와 생애주기를 구분하는 형상관리 전략으로, 저장소 비대화를 막고 배포 재현성을 확보한다. 사용성 문서인 시뮬레이터 매뉴얼은 실제 구동 화면 캡처를 포함하여, 산출물이 최종 사용자와 평가자 모두에게 검증 가능한 형태로 제공되도록 한다~\cite{nielsen1993usability,krug2014dontmakemethink}.

\section{구현을 위한 의사코드}

첫 번째 의사코드는 요구사항 추적성 행렬을 산출물 전반에 걸쳐 구축하는 절차다. 각 요구사항이 설계·구현·시험 산출물의 어느 항목으로 실현되는지를 연결하고, 끊어진 사슬(고아 요구사항, 미검증 항목)을 탐지한다.

\captionof{listing}{의사코드 --- 요구사항 추적성 행렬 구축}
\begin{minted}{text}
추적성_행렬_구축(요구사항들, 산출물들):
    행렬 = 빈_맵()
    각 요구 in 요구사항들:
        행렬[요구.식별자] = { 설계: [], 구현: [], 시험: [] }
    각 산출물 in 산출물들:
        각 항목 in 산출물.항목들:
            각 요구식별자 in 항목.추적태그들:      # 예: REQ-이동-01
                만약 요구식별자 in 행렬:
                    행렬[요구식별자][산출물.종류].추가(항목)
                아니면:
                    경고("미정의 요구 참조", 요구식별자)
    고아 = [r for r in 행렬 if 행렬[r].구현 == 빈배열]
    미검증 = [r for r in 행렬 if 행렬[r].시험 == 빈배열]
    반환 (행렬, 고아, 미검증)     # 끊어진 사슬 보고
\end{minted}
이 절차는 요구사항과 산출물을 양방향으로 연결하여, 구현되지 않은 요구(고아)와 시험되지 않은 요구(미검증)를 명시적 결함으로 드러낸다~\cite{iso29119,mcconnell2004codecomplete}.

두 번째 의사코드는 산출물 형상관리와 품질 게이트를 결합한다. 각 산출물을 버전·상태와 함께 등록하고, 대형 이진 산출물은 릴리스 자산으로 분리하며, 정의된 품질 기준을 통과해야 릴리스로 승격되도록 한다.

\captionof{listing}{의사코드 --- 산출물 형상관리와 품질 게이트}
\begin{minted}{text}
산출물_등록_및_승격(산출물, 저장소, 릴리스):
    산출물.버전 = 다음_버전(저장소, 산출물.식별자)
    만약 산출물.크기 > 대형_임계값:           # 예: 진행보고 PDF
        릴리스.자산_갱신(산출물)              # 저장소 밖 릴리스로 분리
        저장소.기록(산출물.메타데이터만)       # 링크·해시만 커밋
    아니면:
        저장소.기록(산출물)
    통과 = 품질게이트(산출물)                  # 추적성·시험·리뷰 확인
    만약 통과:
        산출물.상태 = "승인"
    아니면:
        산출물.상태 = "보류"; 결함기록(산출물)
    반환 산출물.상태

품질게이트(산출물):
    반환 (산출물.추적성_완결 그리고
           산출물.시험_통과율 >= 기준치 그리고
           산출물.리뷰_승인)
\end{minted}

두 블록은 각각 추적성 검증과 형상·품질 통제를 담당한다. 추적성 행렬이 산출물 간 논리적 완결성을 보장한다면, 형상관리 절차는 산출물의 물리적 저장·버전·품질 상태를 통제하여, 문서 체계 전체가 감사 가능하고 재현 가능하도록 만든다~\cite{chacon2014progit,ieee829}.

\section{기술적·학술적 의의}

첫째, 산출물을 하나의 정합된 체계로 관리한 것은 이 프로젝트가 코드 작성을 넘어 소프트웨어 공학적 성숙도를 지향했음을 보여준다. 코드분석 종합보고서에서 시작해 리팩토링 계획, 머지 점검, 중간 재분석, STR, 매뉴얼, 최종 코드정리, 소스 리스팅으로 이어지는 산출물 흐름은 개발 생애주기의 각 국면을 명시적으로 문서화한 결과다~\cite{mcconnell2004codecomplete,iso25010}. 이러한 문서 체계는 개발자가 바뀌거나 시간이 지나도 프로젝트의 의사결정 맥락이 소실되지 않도록 보존하며, 교육 현장에서 교사·평가자가 시스템을 신뢰하고 채택하는 근거가 된다.

둘째, 요구사항 추적성의 명시화는 이 산출물 체계의 학술적 핵심이다. 미션 요구가 블록·명령 실행기·펌웨어·시뮬레이터·시험으로 이어지는 사슬을 추적 가능하게 만들면, 변경 영향 분석과 회귀 검증이 체계적으로 이뤄진다~\cite{iso29119,fowler2018refactoring}. 특히 모놀리식 시뮬레이터를 모듈로 분해하는 대규모 리팩토링 후에도 요구-구현-시험의 연결이 유지되었음을 중간 재분석과 STR로 확인함으로써, 구조 변경이 기능적 정합성을 훼손하지 않았다는 근거 있는 주장이 가능해진다~\cite{feathers2004legacy}.

셋째, 시험 산출물이 국제 표준과 정렬되었다는 점은 교육용 프로젝트의 품질 주장에 객관성을 부여한다. STR을 ISO/IEC/IEEE 29119·IEEE 829의 문서 틀과 ISO/IEC 25010의 품질 특성에 맞춰 작성함으로써, 품질은 주관적 인상이 아니라 정의된 기준에 대한 측정 결과로 표현된다~\cite{iso29119,ieee829,iso25010}. BLE 프로토콜 처리, 명령 디스패치, 모듈 활성화 같은 핵심 경로를 대상으로 기대-실제 결과를 대조한 시험 근거는, 아동을 대상으로 하는 교구가 갖춰야 할 신뢰성과 안전성의 최소 요건을 문서로 담보한다.

넷째, 형상관리와 산출물 분리 전략은 실무적 지속가능성의 의의를 갖는다. 대형 이진 산출물을 릴리스 자산으로 분리하고 자동화 파이프라인으로 배포하는 방식은 저장소 건전성과 배포 재현성을 동시에 확보하며, 이는 프로젝트가 장기적으로 유지·확장될 때 축적되는 기술 부채를 억제한다~\cite{chacon2014progit,hunt1999pragmatic,ares-repo}. 종합하면 본 보고 기간의 산출물 체계는 코드·문서·시험·형상관리가 요구사항 추적성이라는 축을 중심으로 결합된 하나의 품질 시스템이며, ARES를 일회성 결과물이 아니라 지속 가능한 교육 플랫폼으로 자리매김하게 하는 학술적·기술적 토대다~\cite{iso25010,martin2008cleancode}.


%% ===== 장 심화 확장 — 자동 보강(2차) =====
\section{설계 대안·정량 분석·확장 논의}

산출물 개요 장은 무엇을, 얼마나 상세히 기록할 것인가라는 질문으로 압축된다. \textbf{문서화 수준}에는 명백한 트레이드오프가 존재한다. 과소 문서화는 지식이 특정 개발자의 머릿속에만 남아 인수인계와 유지보수를 마비시키고, 과다 문서화는 코드가 바뀔 때마다 문서가 뒤처져 \emph{거짓을 말하는 문서}가 되어 오히려 신뢰를 무너뜨린다\cite{mcconnell2004codecomplete}. 코드 자체가 최선의 문서라는 원칙\cite{martin2008cleancode}과, 실행 가능한 명세로서의 테스트\cite{beck2003tdd}를 함께 고려하면, 우리가 택한 균형점은 명확하다. \emph{자주 바뀌고 자기설명적인 것은 코드·테스트·타입에 맡기고, 느리게 바뀌며 설계 의도를 담는 것만 산문으로 남긴다.} ARES에서 명령 프로토콜 레퍼런스\cite{ares-repo}, 하드웨어 모듈 패턴, 빌드·배포 절차가 후자에 해당한다. 이는 실용주의 프로그래머\cite{hunt1999pragmatic}가 경계한 \emph{중복(DRY 위반)}—같은 사실을 코드와 문서 두 곳에 적어 동기화 부채를 지는 일—을 피하려는 의도적 선택이다.

\textbf{요구사항 추적성}은 품질 보증의 척추다. ISO/IEC 25010\cite{iso25010}의 품질 특성(기능 적합성·성능·사용성·유지보수성)을 요구사항—설계—구현—시험에 연결하는 추적성 행렬은 특정 요구사항이 어떤 코드와 어떤 시험으로 충족되는지를 증명한다. 그러나 이 행렬은 유지 비용이 만만치 않다. 항목이 $N$개, 산출물이 $M$종이면 추적 링크는 최악의 경우 $N \times M$으로 증가하며, 코드가 한 번 바뀔 때마다 관련 링크를 모두 갱신해야 한다. 소프트웨어 시험 표준\cite{iso29119}과 시험 문서화 규격\cite{ieee829}이 요구하는 완전한 추적성을 초등 교육용 프로젝트에 그대로 적용하면 비용이 편익을 초과한다. 따라서 우리는 \emph{안전·핵심 기능(BLE 명령 파싱, 하드웨어 모듈 활성화)에만 엄밀한 추적을 적용하고 나머지는 완화}하는 위험 기반(risk-based) 접근을 취한다.

\captionof{listing}{의사코드 --- 위험 가중 추적성 우선순위화}
\begin{minted}{text}
함수 추적성_우선순위(요구사항목록):
    작업목록 = 빈_목록
    각 요구사항에 대해 요구사항목록:
        위험 = 심각도(요구사항) * 변경빈도(요구사항)
        만약 위험이 임계값보다 크면:
            작업목록.추가( (요구사항, "완전추적") )   # 핵심 → 코드+시험 링크
        아니면:
            작업목록.추가( (요구사항, "경량추적") )   # 주변 → 커밋 참조만
    위험_내림차순_정렬(작업목록)
    반환 작업목록
\end{minted}

\textbf{형상관리}에서 가장 실무적인 갈등은 대형 이진 산출물의 취급이었다. 진행보고 PDF(약 52MB)나 시뮬 자산 같은 대형 이진 파일을 일반 커밋으로 관리하면, git이 이진 파일의 델타를 효율적으로 저장하지 못해 매 버전이 히스토리를 통째로 부풀린다\cite{chacon2014progit}. 실제로 이 프로젝트에서는 매 커밋마다 52MB PDF가 히스토리에 누적되어 Pages 배포가 6\char`~11분씩 걸리는 문제가 발생했다. 대안은 둘이다. \emph{Git LFS}는 이진 파일을 포인터로 대체해 워킹 트리 경험을 유지하지만 별도 스토리지 할당량과 대역폭 과금, 클라이언트 설치를 요구한다. \emph{릴리스 자산(GitHub Releases)}은 이진물을 버전관리 히스토리 밖으로 완전히 분리해 저장소를 가볍게 유지하되, 파일별 세밀한 버전 이력은 포기한다. 우리는 PDF가 \emph{매 커밋마다 diff를 볼 필요가 없는 최종 산출물}이라는 성격에 근거해 Releases 방식을 택했고, 빌드 스크립트가 고정 태그 자산을 자동 갱신하도록 해 공유 링크를 안정화했다\cite{ares-repo}. 이는 \emph{정본 저장소를 클라우드 동기화 폴더 밖에 두어 손상을 예방}한다는 운영 규칙과 함께, 형상관리를 도구가 아니라 정책으로 다스린 사례다.

이 장의 관리 체계에도 잔여 한계가 있다. 위험 기반 추적성은 초기 위험 판정이 틀리면 정작 중요한 요구사항의 추적을 놓칠 수 있고, 프로젝트가 커지면 완화 대상이 슬그머니 임계값을 넘는데도 갱신되지 않는 \emph{추적 부채}가 쌓인다. Releases 기반 이진 관리는 과거 특정 버전 PDF의 정확한 재현을 어렵게 하므로, 재현성이 중요한 학술 산출물에는 빌드 입력(소스·태그)만 커밋하고 출력을 배제하는 절충이 필요하다. 확장 방향으로는, 레거시 코드 다루기\cite{feathers2004legacy}의 관점에서 문서 없는 기존 코드에 특성화 시험을 붙여 \emph{동작을 문서화}하고, CI 파이프라인이 추적성 링크의 깨짐과 대형 이진물의 우발적 커밋을 자동 검출하도록 계측하는 것이 유효하다. 궁극적으로 문서화·추적성·형상관리는 별개 활동이 아니라, \emph{하나의 변경이 코드·시험·문서·산출물에 어떻게 파급되는지를 일관되게 추적하는 단일 규율}\cite{iso29119}로 수렴해야 한다.

